\section{Limitations and broader impact}

\paragraph{Limitations.} All models are trained on one dataset of mostly frontal, well-lit
YouTube faces. The test sets cover face swaps and reenactment from 2018 to 2020 \todo{and DF40,
if run}; newer diffusion-based swaps and talking-head generators are not tested, and a low score
on them means only that this model did not detect them. Degradations are applied to face crops,
which approximates but does not equal degrading the full video. Calibration fitted on FF++ does
not fully carry over to other data. \todo{add limitations found in the results}

\paragraph{Misuse.} A public detector can be used to tune fakes until they pass, and a low score
can be misrepresented as proof of authenticity. We release the weights for research under a
non-commercial license, with the failure cases and calibration documented, outputs framed as
bands rather than verdicts, and an explicit note that a score is not proof.

\paragraph{Data ethics.} Every dataset is used under its license and only for research; no faces
are redistributed. Many people in these datasets did not consent to appearing in them, which is
one more reason this model must not be used to make claims about specific individuals.

\section{Conclusion}
\todo{Two or three sentences once the results are in: what worked, by how much, and the most
important open problem.}
